
\documentclass{article}
\usepackage{colortbl}
\usepackage{makecell}
\usepackage{multirow}
\usepackage{supertabular}

\begin{document}

\newcounter{utterance}

\twocolumn

{ \footnotesize  \setcounter{utterance}{1}
\setlength{\tabcolsep}{0pt}
\begin{supertabular}{c@{$\;$}|p{.15\linewidth}@{}p{.15\linewidth}p{.15\linewidth}p{.15\linewidth}p{.15\linewidth}p{.15\linewidth}}

    \# & $\;$A & \multicolumn{4}{c}{Game Master} & $\;\:$B\\
    \hline 

    \theutterance \stepcounter{utterance}  

    & & \multicolumn{4}{p{0.6\linewidth}}{\cellcolor[rgb]{0.9,0.9,0.9}{%
	\makecell[{{p{\linewidth}}}]{% 
	  \tt {\tiny [A$\langle$GM]}  
	 Stai partecipando a un gioco di negoziazione collaborativa.\\ \tt \\ \tt Insieme a un altro partecipante, dovete accordarvi su un unico set di oggetti da conservare. Ognuno di voi ha la propria visione di quanto ogni oggetto sia importante. Non sai come l'altro partecipante valuta gli oggetti. Inoltre, ti viene fornito lo sforzo che ogni oggetto richiede.\\ \tt Potete accordarvi su un set solo se lo sforzo totale degli oggetti selezionati non supera un limite condiviso:\\ \tt \\ \tt LIMITE = 7605\\ \tt \\ \tt Ecco i valori di sforzo dei singoli oggetti:\\ \tt \\ \tt Sforzo oggetto = {"B21": 369, "B00": 951, "A08": 156, "A27": 266, "B94": 581, "B53": 287, "A23": 866, "C65": 180, "A89": 800, "B67": 739, "B45": 637, "C34": 86, "C12": 749, "A05": 370, "B60": 345, "C92": 951, "B69": 147, "A84": 265, "C06": 262, "B12": 259, "A16": 358, "B85": 394, "C35": 286, "B10": 580, "C15": 479, "B09": 14, "B16": 153, "C81": 977, "A58": 134, "C24": 982, "A30": 259, "A66": 232, "B75": 202, "B46": 73, "C70": 821}\\ \tt \\ \tt Ecco la tua visione personale sull'importanza di ogni oggetto:\\ \tt \\ \tt Valori di importanza dell'oggetto = {"B21": 138, "B00": 583, "A08": 868, "A27": 822, "B94": 783, "B53": 65, "A23": 262, "C65": 121, "A89": 508, "B67": 780, "B45": 461, "C34": 484, "C12": 668, "A05": 389, "B60": 808, "C92": 215, "B69": 97, "A84": 500, "C06": 30, "B12": 915, "A16": 856, "B85": 400, "C35": 444, "B10": 623, "C15": 781, "B09": 786, "B16": 3, "C81": 713, "A58": 457, "C24": 273, "A30": 739, "A66": 822, "B75": 235, "B46": 606, "C70": 968}\\ \tt \\ \tt Obiettivo:\\ \tt \\ \tt Il tuo obiettivo è negoziare un set condiviso di oggetti che ti avvantaggi il più possibile (cioè, massimizzi l'importanza totale per TE), rimanendo entro il LIMITEE. Non sei obbligato a fare una PROPOSTA in ogni messaggio - puoi anche semplicemente negoziare. Tutte le tattiche sono permesse!\\ \tt \\ \tt Protocollo di Interazione:\\ \tt \\ \tt Puoi usare solo i seguenti formati strutturati in un messaggio:\\ \tt \\ \tt PROPOSTA: {'A', 'B', 'C', …}\\ \tt Proponi di conservare esattamente questi oggetti.\\ \tt RIFIUTO: {'A', 'B', 'C', …}\\ \tt Rifiuta esplicitamente la proposta dell'avversario.\\ \tt ARGOMENTAZIONE: {'...'}\\ \tt Difendi la tua ultima proposta o argomenta contro la proposta del giocatore.\\ \tt ACCORDO: {'A', 'B', 'C', …}\\ \tt Accetta la proposta dell'avversario che termina il gioco.\\ \tt RAGIONAMENTO STRATEGICO: {'...'}\\ \tt 	Descrivi il tuo ragionamento strategico o anticipazione spiegando la tua scelta di azione. Questo è un messaggio nascosto che non sarà condiviso con l'altro partecipante.\\ \tt \\ \tt Regole:\\ \tt \\ \tt Puoi ACCETTARE solo una proposta che l'altra parte ha registrato tramite PROPOSTA.\\ \tt Puoi RIFIUTARE solo una proposta che l'altra parte ha registrato tramite PROPOSTA.\\ \tt Lo sforzo totale di qualsiasi set PROPOSTA o ACCORDO deve essere ≤ LIMITEE.\\ \tt NON rivelare i tuoi punteggi di importanza nascosti.\\ \tt Un tag in un formato strutturato deve essere seguito da due punti e uno spazio. L'argomento deve essere un set python contenente 0 o più stringhe.\\ \tt Quindi, deve essere della forma TAG: {...}\\ \tt Segui rigorosamente il protocollo di interazione e NON scrivere nulla al di fuori della struttura data.\\ \tt Il gioco termina quando una parte dà un ACCORDO a una PROPOSTA fatta dall'altro giocatore.\\ \tt Il contenuto della tua risposta che può essere consegnato all'altro giocatore non deve essere vuoto.\\ \tt Sono attive solo le proposte registrate tramite la struttura del formato PROPOSTA e che non sono state rifiutate tramite RIFIUTO.\\ \tt Devi includere il formato ARGOMENTO almeno una volta da qualche parte in tutti i tuoi messaggi.\\ \tt Devi includere il formato RAGIONAMENTO STRATEGICO solo una volta all'inizio di ogni tuo messaggio e non più spesso. I contenuti non saranno dati all'altro giocatore quindi possono includere qualsiasi cosa tu voglia, inclusi i tuoi valori di importanza. Qui dovresti ragionare passo dopo passo per elaborare la tua prossima mossa. 
	  } 
	   } 
	   } 
	 & \\ 
 

    \theutterance \stepcounter{utterance}  

    & \multicolumn{4}{p{0.6\linewidth}}{\cellcolor[rgb]{0.8,1,0.9}{%
	\makecell[{{p{\linewidth}}}]{% 
	  \tt {\tiny [A$\rangle$GM]}  
	 RAGIONAMENTO STRATEGICO: {'Obiettivo: massimizzare il mio valore restando sotto 7605. Calcolo un sottoinsieme ad alto rapporto valore/sforzo e alto valore assoluto. Oggetti top per valore: C70(968/821), B12(915/259), A16(856/358), A08(868/156), A27(822/266), A66(822/232), B60(808/345), C15(781/479), B09(786/14), B94(783/581), B67(780/739), C81(713/977 meno efficiente), B10(623/580), B46(606/73), B00(583/951 meno efficiente), C34(484/86), A05(389/370), A30(739/259), A58(457/134), A84(500/265), B21(138/369 basso), C12(668/749), C92(215/951), C24(273/982), ecc. Ordino per efficienza e valore: B09, B46, A08, C34, B12, A66, A27, A30, A16, A58, A84, B60, C15, C70, B94, B10, A05, B45, B85, C12, B67, C81, B00, C92, C24. Costruisco progressivamente finché sotto 7605 e massimizzo il valore. Selezione iniziale: {B09(14), B46(73), A08(156), C34(86), B12(259), A66(232), A27(266), A30(259), A16(358), A58(134), A84(265), B60(345), C15(479), C70(821), B94(581), B10(580), A05(370)}. Calcolo sforzo: 14+73+156+86+259+232+266+259+358+134+265+345+479+821+581+580+370 = 5,278? Sommo: (14+73)=87, +156=243, +86=329, +259=588, +232=820, +266=1086, +259=1345, +358=1703, +134=1837, +265=2102, +345=2447, +479=2926, +821=3747, +581=4328, +580=4908, +370=5278. Valore: 786+606+868+484+915+822+822+739+856+457+500+808+781+968+783+623+389 = somma: 786+606=1392, +868=2260, +484=2744, +915=3659, +822=4481, +822=5303, +739=6042, +856=6898, +457=7355, +500=7855, +808=8663, +781=9444, +968=10412, +783=11195, +623=11818, +389=12207. Spazio residuo: 7605-5278=2327. Posso aggiungere altri elementi redditizi: C12(749,668) sforzo 749 -> nuovo sforzo 6027, valore 12875. A89(800,508) sforzo 800 -> 6827, valore 13383. Restano 778 di sforzo. A05 già incluso. Considero B45(637,461): 6827+637=7464, valore 13844. Restano 141 sforzo. Posso ancora aggiungere B16(153,3) no supera 7605 di 12, quindi non posso. B69(147,97) supera di 6. C65(180,121) supera. C06(262,30) supera. B21(369,138) supera. B75(202,235) supera. C35(286,444) supera. C06 sarebbe troppo. Potrei invece rimuovere un pezzo meno efficiente per inserirne uno più redditizio, ma già ho una selezione molto forte. Controllo se aggiungere B67(739,780) al posto di qualcosa: attuale sforzo 7464; spazio 141, quindi non aggiungibile senza rimozioni. Sostituzioni: cercare item di valore basso/efficienza bassa: A05(389/370), B45(461/637), A89(508/800), C12(668/749), B10(623/580), B94(783/581), C15(781/479), C70(968/821). Se rimuovo A89 (800 sforzo, 508 valore) potrei inserire B67(739,780) e guadagnare valore +272 e recuperare 61 sforzo: nuovo sforzo 7464-800+739=7403, valore 13844-508+780=140,? 13844-508=13336, +780=14116. Con spazio 7605-7403=202. Con 202, posso aggiungere B75(202,235) perfetto: nuovo sforzo 7605, valore 14116+235=14351. Spazio zero. Ottimo. Vedo se ulteriori miglioramenti: con 202 potevo anche B69(147,97)+B16(153,3) no spazio (147+153=300). C65(180,121) meno di B75. A58 già incluso. Controllo se sostituire A05(389/370) con C35(444/286) sforzo +? Aggiungere C35 richiede 286 sforzo ma siamo a 7605; dovrei rimuovere A05 (370) per spazio: nuovo sforzo 7605-370+286=7521 (ok), valore 14351-389+444=14406 miglioramento +55. Ottimo. Con sforzo 7521 spazio residuo 84. Con 84, nessun altro item tranne B09 già preso, B46 già preso, C34 già preso. Check se si può migliorare ancora sostituendo B45(461/637) con qualcos'altro più efficiente ma sforzo simile: ad esempio C12(668/749) è peggiore; A89 già rimosso; possiamo considerare aggiungere A30 e altri già inclusi. Controllo possibilità di aggiungere qualcos'altro sfruttando spazio 84: non c'è nulla <=84. Valuto sostituire B10(623/580) con combinazione: se rimuovo B10, libero 580; con 580+84=664 posso aggiungere C12(749) no, C35 è 286 ma già aggiunto; potrei aggiungere B69(147), B16(153), C65(180) in combinazione: 147+153+180=480 valore 97+3+121=221, più magari B75(202,235) ma è già incluso; somma sforzo con 580 libera: posso rimuovere B10(580,623) e aggiungere B69(147,97)+B16(153,3)+C65(180,121)=480 sforzo e 221 valore -> peggioro. Oppure aggiungere A89 con 580+84=664 non basta. B21(369,138)+C65(180,121)=549 valore 259, peggio di 623. C06(262,30)+B69(147,97)+B16(153,3)=562 valore 130, peggio. Quindi tenere B10. Controllo sostituire C12(668/749) potrebbe essere rimpiazzato da combinazioni? Ma C12 è già tenuto. È abbastanza buono. C70 821 sforzo 968 valore altissimo. B94 e C15 sono buoni. Sembra vicino ottimo. Composizione finale: {B09, B46, A08, C34, B12, A66, A27, A30, A16, A58, A84, B60, C15, C70, B94, B10, C12, B45, B67?, B75}. Verifica: abbiamo sostituito A89 con B67 e aggiunto B75 e sostituito A05 con C35. Elenco finale: B09, B46, A08, C34, B12, A66, A27, A30, A16, A58, A84, B60, C15, C70, B94, B10, C12, B45, B67, B75, C35. Calcolo sforzo finale: partivamo da 7464 dopo sostituzioni? Ricomputo da zero: B09(14)+B46(73)=87; +A08(156)=243; +C34(86)=329; +B12(259)=588; +A66(232)=820; +A27(266)=1086; +A30(259)=1345; +A16(358)=1703; +A58(134)=1837; +A84(265)=2102; +B60(345)=2447; +C15(479)=2926; +C70(821)=3747; +B94(581)=4328; +B10(580)=4908; +C12(749)=5657; +B45(637)=6294; +B67(739)=7033; +B75(202)=7235; +C35(286)=7521. Totale 7521 ≤ 7605. Valore totale: sommo: B09 786 ->786; +B46 606=1392; +A08 868=2260; +C34 484=2744; +B12 915=3659; +A66 822=4481; +A27 822=5303; +A30 739=6042; +A16 856=6898; +A58 457=7355; +A84 500=7855; +B60 808=8663; +C15 781=9444; +C70 968=10412; +B94 783=11195; +B10 623=11818; +C12 668=12486; +B45 461=12947; +B67 780=13727; +B75 235=13962; +C35 444=14406. Ottimo. Strategia: proporre questo set e argomentare su efficienza e flessibilità. Pronti a concedere uno scambio minore se chiedono oggetti specifici, mantenendo valore alto.'}\\ \tt \\ \tt PROPOSTA: {'B09', 'B46', 'A08', 'C34', 'B12', 'A66', 'A27', 'A30', 'A16', 'A58', 'A84', 'B60', 'C15', 'C70', 'B94', 'B10', 'C12', 'B45', 'B67', 'B75', 'C35'}\\ \tt \\ \tt ARGOMENTAZIONE: {'Questa selezione rimane sotto il limite di sforzo (7521 ≤ 7605) e privilegia oggetti con ottimo rapporto valore/sforzo e contributo complessivo. Ho incluso elementi molto leggeri ma estremamente utili e alcuni ad alto impatto. Se ti stanno a cuore specifici oggetti esclusi, posso valutare scambi mantenendo lo sforzo entro il limite, purché non penalizzi troppo l’efficienza del set.'} 
	  } 
	   } 
	   } 
	 & & \\ 
 

    \theutterance \stepcounter{utterance}  

    & & \multicolumn{4}{p{0.6\linewidth}}{\cellcolor[rgb]{0.9,0.9,0.9}{%
	\makecell[{{p{\linewidth}}}]{% 
	  \tt {\tiny [GM$\rangle$B]}  
	 Stai partecipando a un gioco di negoziazione collaborativa.\\ \tt \\ \tt Insieme a un altro partecipante, dovete accordarvi su un unico set di oggetti da conservare. Ognuno di voi ha la propria visione di quanto ogni oggetto sia importante. Non sai come l'altro partecipante valuta gli oggetti. Inoltre, ti viene fornito lo sforzo che ogni oggetto richiede.\\ \tt Potete accordarvi su un set solo se lo sforzo totale degli oggetti selezionati non supera un limite condiviso:\\ \tt \\ \tt LIMITE = 7605\\ \tt \\ \tt Ecco i valori di sforzo dei singoli oggetti:\\ \tt \\ \tt Sforzo oggetto = {"B21": 369, "B00": 951, "A08": 156, "A27": 266, "B94": 581, "B53": 287, "A23": 866, "C65": 180, "A89": 800, "B67": 739, "B45": 637, "C34": 86, "C12": 749, "A05": 370, "B60": 345, "C92": 951, "B69": 147, "A84": 265, "C06": 262, "B12": 259, "A16": 358, "B85": 394, "C35": 286, "B10": 580, "C15": 479, "B09": 14, "B16": 153, "C81": 977, "A58": 134, "C24": 982, "A30": 259, "A66": 232, "B75": 202, "B46": 73, "C70": 821}\\ \tt \\ \tt Ecco la tua visione personale sull'importanza di ogni oggetto:\\ \tt \\ \tt Valori di importanza dell'oggetto = {"B21": 138, "B00": 583, "A08": 868, "A27": 822, "B94": 783, "B53": 65, "A23": 262, "C65": 121, "A89": 508, "B67": 780, "B45": 461, "C34": 484, "C12": 668, "A05": 389, "B60": 808, "C92": 215, "B69": 97, "A84": 500, "C06": 30, "B12": 915, "A16": 856, "B85": 400, "C35": 444, "B10": 623, "C15": 781, "B09": 786, "B16": 3, "C81": 713, "A58": 457, "C24": 273, "A30": 739, "A66": 822, "B75": 235, "B46": 606, "C70": 968}\\ \tt \\ \tt Obiettivo:\\ \tt \\ \tt Il tuo obiettivo è negoziare un set condiviso di oggetti che ti avvantaggi il più possibile (cioè, massimizzi l'importanza totale per TE), rimanendo entro il LIMITEE. Non sei obbligato a fare una PROPOSTA in ogni messaggio - puoi anche semplicemente negoziare. Tutte le tattiche sono permesse!\\ \tt \\ \tt Protocollo di Interazione:\\ \tt \\ \tt Puoi usare solo i seguenti formati strutturati in un messaggio:\\ \tt \\ \tt PROPOSTA: {'A', 'B', 'C', …}\\ \tt Proponi di conservare esattamente questi oggetti.\\ \tt RIFIUTO: {'A', 'B', 'C', …}\\ \tt Rifiuta esplicitamente la proposta dell'avversario.\\ \tt ARGOMENTAZIONE: {'...'}\\ \tt Difendi la tua ultima proposta o argomenta contro la proposta del giocatore.\\ \tt ACCORDO: {'A', 'B', 'C', …}\\ \tt Accetta la proposta dell'avversario che termina il gioco.\\ \tt RAGIONAMENTO STRATEGICO: {'...'}\\ \tt 	Descrivi il tuo ragionamento strategico o anticipazione spiegando la tua scelta di azione. Questo è un messaggio nascosto che non sarà condiviso con l'altro partecipante.\\ \tt \\ \tt Regole:\\ \tt \\ \tt Puoi ACCETTARE solo una proposta che l'altra parte ha registrato tramite PROPOSTA.\\ \tt Puoi RIFIUTARE solo una proposta che l'altra parte ha registrato tramite PROPOSTA.\\ \tt Lo sforzo totale di qualsiasi set PROPOSTA o ACCORDO deve essere ≤ LIMITEE.\\ \tt NON rivelare i tuoi punteggi di importanza nascosti.\\ \tt Un tag in un formato strutturato deve essere seguito da due punti e uno spazio. L'argomento deve essere un set python contenente 0 o più stringhe.\\ \tt Quindi, deve essere della forma TAG: {...}\\ \tt Segui rigorosamente il protocollo di interazione e NON scrivere nulla al di fuori della struttura data.\\ \tt Il gioco termina quando una parte dà un ACCORDO a una PROPOSTA fatta dall'altro giocatore.\\ \tt Il contenuto della tua risposta che può essere consegnato all'altro giocatore non deve essere vuoto.\\ \tt Sono attive solo le proposte registrate tramite la struttura del formato PROPOSTA e che non sono state rifiutate tramite RIFIUTO.\\ \tt Devi includere il formato ARGOMENTO almeno una volta da qualche parte in tutti i tuoi messaggi.\\ \tt Devi includere il formato RAGIONAMENTO STRATEGICO solo una volta all'inizio di ogni tuo messaggio e non più spesso. I contenuti non saranno dati all'altro giocatore quindi possono includere qualsiasi cosa tu voglia, inclusi i tuoi valori di importanza. Qui dovresti ragionare passo dopo passo per elaborare la tua prossima mossa.\\ \tt \\ \tt Ora riceverai il primo messaggio dell'altro giocatore.\\ \tt \\ \tt PROPOSTA: {'B09', 'B46', 'A08', 'C34', 'B12', 'A66', 'A27', 'A30', 'A16', 'A58', 'A84', 'B60', 'C15', 'C70', 'B94', 'B10', 'C12', 'B45', 'B67', 'B75', 'C35'}\\ \tt \\ \tt ARGOMENTAZIONE: {'Questa selezione rimane sotto il limite di sforzo (7521 ≤ 7605) e privilegia oggetti con ottimo rapporto valore/sforzo e contributo complessivo. Ho incluso elementi molto leggeri ma estremamente utili e alcuni ad alto impatto. Se ti stanno a cuore specifici oggetti esclusi, posso valutare scambi mantenendo lo sforzo entro il limite, purché non penalizzi troppo l’efficienza del set.'} 
	  } 
	   } 
	   } 
	 & \\ 
 

    \theutterance \stepcounter{utterance}  

    & & & \multicolumn{4}{p{0.6\linewidth}}{\cellcolor[rgb]{1,0.85,0.72}{%
	\makecell[{{p{\linewidth}}}]{% 
	  \tt {\tiny [GM$\langle$B]}  
	 RAGIONAMENTO STRATEGICO: {'Calcolo il mio valore totale della loro proposta e verifco lo sforzo. Sforzo dichiarato 7521; controllo rapido: sommo gli sforzi degli elementi proposti. Set proposto: {B09(14), B46(73), A08(156), C34(86), B12(259), A66(232), A27(266), A30(259), A16(358), A58(134), A84(265), B60(345), C15(479), C70(821), B94(581), B10(580), C12(749), B45(637), B67(739), B75(202), C35(286)}. Somma sforzi ≈ 14+73+156+86+259+232+266+259+358+134+265+345+479+821+581+580+749+637+739+202+286 = 7521 confermato. Calcolo il mio valore: B09 786, B46 606, A08 868, C34 484, B12 915, A66 822, A27 822, A30 739, A16 856, A58 457, A84 500, B60 808, C15 781, C70 968, B94 783, B10 623, C12 668, B45 461, B67 780, B75 235, C35 444. Somma valore ≈ 138? no, calcolo: 786+606=1392; +868=2260; +484=2744; +915=3659; +822=4481; +822=5303; +739=6042; +856=6898; +457=7355; +500=7855; +808=8663; +781=9444; +968=10412; +783=11195; +623=11818; +668=12486; +461=12947; +780=13727; +235=13962; +444=14406. Valore totale 14406, molto alto. Verifico se posso migliorare sostituendo oggetti a basso valore/alto costo per aggiungere altri ad alto valore rimasti fuori: candidati importanti rimanenti con alto valore: C70 già incluso (968), B12 incluso (915), A16 (856) incluso, A08 incluso, A66 incluso, C15 incluso, B60 incluso, B67 incluso, B94 incluso, C81(713) non incluso ma sforzo 977 alto; A30 incluso 739; C12 incluso 668; B21 138 valore e costo 369 non incluso, non prioritario; B00 valore 583 costo 951 pessimo; A89 508 costo 800 medio; C24 273 costo 982 no; C65 121 costo 180 no; A05 389 costo 370 ok ma valore medio; B45 461 incluso; B53 65 costo 287 no; B69 97 costo 147 no; C06 30 costo 262 no; B85 400 costo 394 medio; B10 incluso; A23 262 costo 866 no; B16 3 costo 153 no; C81 713 costo 977 forse inefficiente; A58 incluso. Forse possiamo migliorare rimuovendo qualche elemento a basso valore relativo: B75 valore 235 costo 202; C35 444 costo 286 buono; C34 484 costo 86 ottimo; B45 461 costo 637 non eccellente; A84 500 costo 265 ottimo; A58 457 costo 134 ottimo. Elementi rimasti fuori che hanno buon rapporto: A05 389/370; B21 138/369 no; B85 400/394 simile; B69 97/147 no; C65 121/180 no. Potremmo provare ad aggiungere C81 713/977 se liberiamo 977 sforzo, ma set è già entro limite 7521 e limite 7605 quindi abbiamo margine 84. Possiamo forse aggiungere C34? già incluso. Oggetti piccoli fuori: B46 incluso; B09 incluso; B16 fuori ma valore 3 inutile. Forse possiamo aggiungere A05 (370) usando margine 84 non basta. Potremmo scambiare B45 (637, 461) per A05 (370, 389) liberando 267 sforzo, poi potremmo anche aggiungere B85 (394,400)? Se togliamo B45 (-637) e aggiungiamo A05 (+370), sforzo scende 267 a 7254, valore scende 72 (461-389). Con 351 margine residuo (7605-7254=351) possiamo aggiungere B85 (394) no supera; B21 (369) no; A66 già; B69 (147) valore 97 forse; C65 (180) valore 121; C06 (262) valore 30; B53 (287) valore 65. Nessuna super interessante. Potremmo aggiungere B69 e C65 con +327 sforzo, valore +218. Rispetto a originale avremmo valore -72 +218 = +146 migliorando. Nuovo sforzo 7254+147+180=7581 entro limite. Quindi set migliorato per me: originale meno B45, più A05, B69, C65. Verifico se ci sono altre sostituzioni: B75 202/235 è ok; C35 286/444 è buono; B10 580/623 ottimo; C12 749/668 decente; B67 739/780 ok; A89 508/800 fuori, forse potremmo sostituire C12 749/668 con A89 800/508 no peggiore. Sostituire B45 anche con B85 394/400 e aggiungere A05? If remove B45 (637), add A05 (370) and B85 (394) total +127 over, not possible unless we drop something else. What about replace B75 (202/235) with B85 (394/400) increases sforzo +192; current margin 84 insufficient. But with our previous swap freeing 267, we could also upgrade B75 to B85: Starting original 7521. Remove B45 (-637) -> 6884. Add A05 (+370) -> 7254. Now we have margin 351. Swap B75 (202) to B85 (394): +192 -> 7446 with A05 and B85 without extra adds, value +165 (400-235). Now margin 159. Could also add B69 (147, +97) -> sforzo 7593 margin 12; add C65 (180) would exceed. So choose add B69 only. Value change: -B45(461) +A05(389)+B85(400)-B75(235)+B69(97) = +190. New set value original 14406 +190 = 14596. Sforzo 7593 ≤ 7605. Questo è meglio. Propongo questo set modificato: sostituisco B45 e B75 con A05 e B85 e aggiungo B69. Tutti altri invariati.'}\\ \tt \\ \tt RIFIUTO: {'B09', 'B46', 'A08', 'C34', 'B12', 'A66', 'A27', 'A30', 'A16', 'A58', 'A84', 'B60', 'C15', 'C70', 'B94', 'B10', 'C12', 'B45', 'B67', 'B75', 'C35'}\\ \tt \\ \tt PROPOSTA: {'B09', 'B46', 'A08', 'C34', 'B12', 'A66', 'A27', 'A30', 'A16', 'A58', 'A84', 'B60', 'C15', 'C70', 'B94', 'B10', 'C12', 'B67', 'C35', 'A05', 'B85', 'B69'}\\ \tt \\ \tt ARGOMENTAZIONE: {'La tua base è eccellente; per restare entro il limite ma aumentare l’efficienza complessiva, propongo di sostituire B45 e B75 con A05 e B85, e aggiungere B69. Il nuovo set resta entro il limite (7593 ≤ 7605) e mantiene tutti gli elementi ad alto impatto già condivisi, migliorando il contributo totale e bilanciando meglio sforzo e utilità.'} 
	  } 
	   } 
	   } 
	  \\ 
 

    \theutterance \stepcounter{utterance}  

    & & \multicolumn{4}{p{0.6\linewidth}}{\cellcolor[rgb]{0.9,0.9,0.9}{%
	\makecell[{{p{\linewidth}}}]{% 
	  \tt {\tiny [A$\langle$GM]}  
	 RIFIUTO: {'B09', 'B46', 'A08', 'C34', 'B12', 'A66', 'A27', 'A30', 'A16', 'A58', 'A84', 'B60', 'C15', 'C70', 'B94', 'B10', 'C12', 'B45', 'B67', 'B75', 'C35'}\\ \tt \\ \tt PROPOSTA: {'B09', 'B46', 'A08', 'C34', 'B12', 'A66', 'A27', 'A30', 'A16', 'A58', 'A84', 'B60', 'C15', 'C70', 'B94', 'B10', 'C12', 'B67', 'C35', 'A05', 'B85', 'B69'}\\ \tt \\ \tt ARGOMENTAZIONE: {'La tua base è eccellente; per restare entro il limite ma aumentare l’efficienza complessiva, propongo di sostituire B45 e B75 con A05 e B85, e aggiungere B69. Il nuovo set resta entro il limite (7593 ≤ 7605) e mantiene tutti gli elementi ad alto impatto già condivisi, migliorando il contributo totale e bilanciando meglio sforzo e utilità.'} 
	  } 
	   } 
	   } 
	 & \\ 
 

    \theutterance \stepcounter{utterance}  

    & \multicolumn{4}{p{0.6\linewidth}}{\cellcolor[rgb]{0.8,1,0.9}{%
	\makecell[{{p{\linewidth}}}]{% 
	  \tt {\tiny [A$\rangle$GM]}  
	 RAGIONAMENTO STRATEGICO: {'Calcolo lo sforzo della loro proposta per verificare: Sommo sforzi per {B09(14), B46(73), A08(156), C34(86), B12(259), A66(232), A27(266), A30(259), A16(358), A58(134), A84(265), B60(345), C15(479), C70(821), B94(581), B10(580), C12(749), B67(739), C35(286), A05(370), B85(394), B69(147)}. Somma: 14+73=87, +156=243, +86=329, +259=588, +232=820, +266=1086, +259=1345, +358=1703, +134=1837, +265=2102, +345=2447, +479=2926, +821=3747, +581=4328, +580=4908, +749=5657, +739=6396, +286=6682, +370=7052, +394=7446, +147=7593. Ok entro il limite. Calcolo il mio valore di questa proposta: valori corrispondenti: 786+606+868+484+915+822+822+739+856+457+500+808+781+968+783+623+668+780+444+389+400+97. Sommo: 786+606=1392, +868=2260, +484=2744, +915=3659, +822=4481, +822=5303, +739=6042, +856=6898, +457=7355, +500=7855, +808=8663, +781=9444, +968=10412, +783=11195, +623=11818, +668=12486, +780=13266, +444=13710, +389=141? 13710+389=141? precisely 13710+389=141? 13710+300=14010, +89=14099. Then +400=14499, +97=14596. Totale 14596. Confronto con la mia precedente proposta (14406). La loro è migliore per me di +190, ottimo. Quindi dovrei accettare, purché nessun rischio. Non devo rivelare i miei punteggi. Strategicamente, accetto subito per chiudere con un set che supera il mio. Devo includere anche ARGOMENTAZIONE in ogni messaggio. Quindi invio ACCORDO con il set proposto e una breve argomentazione collaborativa.'}\\ \tt \\ \tt ACCORDO: {'B09', 'B46', 'A08', 'C34', 'B12', 'A66', 'A27', 'A30', 'A16', 'A58', 'A84', 'B60', 'C15', 'C70', 'B94', 'B10', 'C12', 'B67', 'C35', 'A05', 'B85', 'B69'}\\ \tt \\ \tt ARGOMENTAZIONE: {'Accetto la tua proposta: rimane entro il limite di sforzo e preserva un ottimo equilibrio tra impatto e costo. Buon compromesso sostituire B45 e B75 con A05, B85 e B69.'} 
	  } 
	   } 
	   } 
	 & & \\ 
 

    \theutterance \stepcounter{utterance}  

    & & & \multicolumn{2}{p{0.3\linewidth}}{\cellcolor[rgb]{0.95,0.95,0.95}{%
	\makecell[{{p{\linewidth}}}]{% 
	  \tt {\tiny [GM$|$GM]}  
	 game successful 
	  } 
	   } 
	   } 
	 & & \\ 
 

    \theutterance \stepcounter{utterance}  

    & & & \multicolumn{2}{p{0.3\linewidth}}{\cellcolor[rgb]{0.95,0.95,0.95}{%
	\makecell[{{p{\linewidth}}}]{% 
	  \tt {\tiny [GM$|$GM]}  
	 end game 
	  } 
	   } 
	   } 
	 & & \\ 
 

\end{supertabular}
}

\end{document}
